\chapter{Cơ sở lý thuyết}
\label{chap:co-so-ly-thuyet}

\section{Tổng quan quản lý và cấu hình thiết bị mạng}

Quản lý thiết bị mạng bao gồm lưu trữ thông tin thiết bị, theo dõi trạng thái, thu thập cấu hình và thực hiện thay đổi cấu hình khi cần.

\section{Giao thức và kỹ thuật liên quan}

\subsection{SSH và Telnet}

SSH là giao thức truy cập từ xa có mã hóa, phù hợp hơn Telnet trong môi trường thực tế.

\subsection{VLAN, DHCP, định tuyến và ACL}

Các chức năng cấu hình mạng thường gặp gồm VLAN, DHCP, định tuyến tĩnh, định tuyến động, NAT và ACL.

\section{Python trong tự động hóa mạng}

Python thường được dùng trong tự động hóa mạng nhờ cú pháp đơn giản, hệ sinh thái thư viện phong phú và khả năng tích hợp với cơ sở dữ liệu, giao diện và API.

\section{Thư viện sử dụng}

Có thể sử dụng các thư viện như Paramiko, Netmiko, PyQt6, SQLite hoặc SQL Server tùy theo phạm vi phần mềm.
